\documentclass{article}

\PassOptionsToPackage{numbers}{natbib}
\usepackage[preprint]{neurips_2026}

\usepackage[utf8]{inputenc}
\usepackage[T1]{fontenc}
\usepackage{hyperref}
\usepackage{url}
\usepackage{booktabs}
\usepackage{amsfonts}
\usepackage{amsmath}
\usepackage{amssymb}
\usepackage{nicefrac}
\usepackage{microtype}
\usepackage{xcolor}
\usepackage{graphicx}
\usepackage{enumitem}
\usepackage{subcaption}
\usepackage{multirow}

\graphicspath{{figures/}}

\title{Experimental Results: PIGNN-Attn-LS across Power System Cases}

\author{Changhun Kim}

\begin{document}
\maketitle

%==============================================================================
\section{Experimental Setup}
%==============================================================================

\subsection{Baseline Definition}

The \textbf{baseline} throughout all experiments (case14--case300) is the
\emph{DC power-flow initialisation} stored as \texttt{u\_start} in each
parquet file.  Concretely, voltage magnitudes are set to per-bus nominal
values (near $1.0$\,p.u.) and voltage angles are taken from the DC
power-flow solution of the same operating scenario.  The parquet filenames
confirm this: they all carry the suffix \texttt{dc\_compile\_siNR}.
The baseline RMSE values reported below are the errors of this DC-initialised
guess against the Newton--Raphson converged solution \emph{before} the model
does any work.  The model's job is to reduce these errors.

Note that the LVN Heo1 experiments used a genuine \emph{flat start}
($|V|=1.0$\,p.u., $\theta=0^\circ$) precisely because reliable DC
initialisation was not available for that network.

\subsection{Test Cases: Bus and Branch Counts}

Table~\ref{tab:grid_size} lists the $N_{\mathrm{bus}}$ and $N_{\mathrm{branch}}$
for each test case.  The branch count determines the number of edges in the
graph and directly affects the cost of the edge-aware attention.

\begin{table}[h]
\centering
\caption{Grid size of test cases.}
\label{tab:grid_size}
\begin{tabular}{lrrl}
\toprule
Case & $N_{\mathrm{bus}}$ & $N_{\mathrm{branch}}$ & Notes \\
\midrule
case14\_ppcY          & 14  & 20  & IEEE 14-bus \\
case24\_ieee\_rts     & 24  & 38  & IEEE RTS-24 \\
case30\_ppcY\_A       & 30  & 41  & IEEE 30-bus \\
case33bw\_ppcY\_A     & 33  & 32  & Radial distribution \\
case39\_ppcY\_A       & 39  & 46  & New England 39-bus \\
case57\_ppcY\_A       & 57  & 80  & IEEE 57-bus \\
case89pegase          & 89  & 210 & PEGASE 89-bus (dense) \\
case118\_ppcY\_A      & 118 & 186 & IEEE 118-bus \\
case145\_ppcY\_A      & 145 & 453 & PEGASE 145-bus (very dense) \\
case300\_ppcY\_A      & 300 & 411 & PEGASE 300-bus (run on A100) \\
LVN Heo1              & 722 & 796 & LEW Verteilnetz GmbH distribution network (reduced internal PPC) \\
\bottomrule
\end{tabular}
\end{table}

\subsubsection*{LVN Heo1 Network Topology: From CGMES to Internal PPC}
\label{subsec:lvn_topology}

The LVN Heo1 grid originates as a CGMES XML file and passes through a
multi-stage reduction before reaching the Newton--Raphson / PIGNN solver.
Table~\ref{tab:lvn_cgmes} summarises the raw CGMES object counts,
Table~\ref{tab:lvn_pandapower} the pandapower equipment counts after import,
and Table~\ref{tab:lvn_ppc} the final internal PPC system used for power
flow.

\begin{table}[h]
\centering
\caption{LVN Heo1 -- raw CGMES object counts.}
\label{tab:lvn_cgmes}
\begin{tabular}{lr}
\toprule
CGMES Object Type & Count \\
\midrule
ConnectivityNode  & 3{,}782 \\
TopologicalNode   & 1{,}197 \\
Terminal          & 9{,}967 \\
BusbarSection     & 682 \\
Breaker           & 895 \\
Disconnector      & 184 \\
PowerTransformer  & 216 \\
ACLineSegment     & 608 \\
EquivalentBranch  & 22 \\
EnergyConsumer    & 275 \\
\bottomrule
\end{tabular}
\end{table}

\begin{table}[h]
\centering
\caption{LVN Heo1 -- pandapower network after CGMES import and Model-A
  cleanup.}
\label{tab:lvn_pandapower}
\setlength{\tabcolsep}{5pt}
\begin{tabular}{lrr}
\toprule
Element & Total & In-service \\
\midrule
\multicolumn{3}{l}{\textit{Buses and switches}} \\
\quad net.bus (ConnectivityNode)  & 3{,}782 & 3{,}200 \\
\quad bus--bus switches (all)     & 3{,}630 & --- \\
\quad bus--bus switches (closed)  & 2{,}600 & --- \\
\midrule
\multicolumn{3}{l}{\textit{Branches}} \\
\quad ACLineSegment (line)        & 608 & 606 \\
\quad 2-winding transformer       & 200 & --- \\
\quad 3-winding transformer       & 16  & --- \\
\quad Impedance element           & 22  & --- \\
\bottomrule
\end{tabular}
\end{table}

\begin{table}[h]
\centering
\caption{LVN Heo1 -- internal PPC system used for Newton--Raphson and PIGNN.
  The large bus reduction ($3{,}782 \to 722$) is due to closed bus--bus
  switches fusing many ConnectivityNodes into single electrical buses.
  The branch count ($796$) reflects only active, in-service elements in the
  solved connected component.}
\label{tab:lvn_ppc}
\setlength{\tabcolsep}{5pt}
\begin{tabular}{lrrl}
\toprule
Stage & Buses & Branches & Notes \\
\midrule
External PPC (pandapower compile) & 3{,}798 & 878 &
  $3{,}782 + 16$ aux buses for 3W-trafos \\
Unique fused bus groups           & 1{,}304 & ---  &
  after closed-switch merging \\
\textbf{Internal PPC (NR / PIGNN)} & \textbf{722} & \textbf{796} &
  active buses/branches only \\
\midrule
\multicolumn{4}{l}{\textit{Internal branch breakdown}} \\
\quad Line (ACLineSegment)        & --- & 581 & \\
\quad 2-winding transformer       & --- & 145 & \\
\quad 3-winding transformer       & --- & 48  & $16\times3$ legs \\
\quad Impedance element           & --- & 22  & \\
\bottomrule
\end{tabular}
\end{table}

%==============================================================================
\section{Full Baseline vs.\ Model Comparison}
%==============================================================================

Table~\ref{tab:baseline_vs_model} summarises the RMSE of the DC-initialised
baseline and the trained model across all cases.  All model runs use
$K=40$, $d=4$, $d_{\mathrm{hi}}=16$, \texttt{numattn}$=1$, Armijo$=\mathrm{True}$,
100 epochs, train ratio $= 1/3$.

\begin{table}[h]
\centering
\small
\caption{DC-initialised baseline vs.\ PIGNN-Attn-LS model.  Reductions are
  $\lfloor$baseline\,/\,model$\rfloor$.}
\label{tab:baseline_vs_model}
\setlength{\tabcolsep}{4pt}
\begin{tabular}{lrrrcrrc}
\toprule
\multirow{2}{*}{Case} & \multirow{2}{*}{$N_{\mathrm{bus}}$}
  & \multicolumn{2}{c}{$|V|$ RMSE (p.u.)} & \multirow{2}{*}{Reduc.}
  & \multicolumn{2}{c}{$\theta$ RMSE (${}^\circ$)} & \multirow{2}{*}{Reduc.} \\
\cmidrule(lr){3-4}\cmidrule(lr){6-7}
& & Baseline & Model & & Baseline & Model & \\
\midrule
case24\_ieee\_rts  & 24  & $2.24\times10^{-2}$ & $1\times10^{-5}$ & $\sim2000\times$ & 0.91  & 0.01 & $\sim90\times$ \\
case30\_ppcY\_A    & 30  & $1.83\times10^{-2}$ & $7\times10^{-5}$ & $\sim260\times$  & 0.21  & 0.025 & $\sim8\times$ \\
case33bw\_ppcY\_A  & 33  & $5.96\times10^{-2}$ & $5\times10^{-4}$ & $\sim120\times$  & 1.47  & 0.10  & $\sim15\times$ \\
case39\_ppcY\_A    & 39  & $1.41\times10^{-2}$ & $5\times10^{-5}$ & $\sim280\times$  & 1.50  & 0.07  & $\sim21\times$ \\
case57\_ppcY\_A    & 57  & $1.26\times10^{-1}$ & $3\times10^{-3}$ & $\sim42\times$   & 1.13  & 0.30  & $\sim4\times$ \\
case89pegase       & 89  & $2.03\times10^{-2}$ & $2.5\times10^{-3}$ & $\sim8\times$  & 0.30  & 0.27  & $\sim1.1\times$ \\
case118\_ppcY\_A   & 118 & $9.02\times10^{-3}$ & $6\times10^{-5}$ & $\sim150\times$  & 2.73  & 0.08  & $\sim34\times$ \\
case145\_ppcY\_A   & 145 & $6.52\times10^{-2}$ & $1.77\times10^{-3}$ & $\sim37\times$ & 18.64 & 1.77  & $\sim11\times$ \\
case300\_ppcY\_A   & 300 & $4.82\times10^{-2}$ & $9.80\times10^{-3}$ & $\sim5\times$ & 19.58 & 1.77  & $\sim11\times$ \\
LVN Heo1          & 722 & $4.03\times10^{-2}$ & \multicolumn{2}{c}{---\,(failed$^\dagger$)} & 3.44  & \multicolumn{2}{c}{---\,(failed$^\dagger$)} \\
\bottomrule
\multicolumn{8}{l}{${}^\dagger$ All LVN training runs failed to learn: no-Armijo runs diverged to $\theta\!\approx\!91^\circ$;}\\
\multicolumn{8}{l}{\phantom{${}^\dagger$} Armijo runs had all steps rejected (physics loss never decreased). See Section~\ref{sec:lvn}.}\\
\end{tabular}
\end{table}

\paragraph{Reading the baseline.}
The DC-initialised error is \emph{not} zero because the DC power-flow is an
approximation: it ignores reactive power, line losses, and voltage magnitude
deviations, so the DC angles and near-nominal magnitudes do not exactly match
the AC Newton--Raphson solution.  The residual grows with grid complexity and
loading: case145 retains a DC-initialised angle error of $18.64^\circ$ because
the PEGASE 145-bus network has strong reactive-power variation and large angle
spreads.  case57 has a large DC-initialised $|V|$ error ($0.126$\,p.u.)
because several buses operate at voltages well below nominal at the AC
solution.

%==============================================================================
\section{Armijo vs.\ Non-Armijo}
%==============================================================================

The backtracking Armijo line search selects a step size $\alpha \in
\{1, \rho, \rho^2, \ldots\}$ at each unrolled step that satisfies the
sufficient-decrease condition on the power-mismatch $\ell^\infty$-norm:
\[
  F(x + \alpha\,\Delta x) \;\le\; (1 - c_1\,\alpha)\,F(x).
\]
Without Armijo (\texttt{armijo=False}), the full proposed step $\alpha=1$
is always accepted, which can cause oscillation or divergence when the
update overshoots.  Table~\ref{tab:armijo} shows the effect across cases.

\begin{table}[h]
\centering
\caption{Effect of Armijo line search. $K=40$, $d=4$, $d_{\mathrm{hi}}=16$,
  \texttt{numattn}$=1$, 100\,ep.}
\label{tab:armijo}
\begin{tabular}{lrrrrl}
\toprule
\multirow{2}{*}{Case} & \multicolumn{2}{c}{Without Armijo} & \multicolumn{2}{c}{With Armijo} & \multirow{2}{*}{Verdict} \\
\cmidrule(lr){2-3}\cmidrule(lr){4-5}
& $|V|$ RMSE & $\theta$ (${}^\circ$) & $|V|$ RMSE & $\theta$ (${}^\circ$) & \\
\midrule
case24\_ieee\_rts & $6\times10^{-4}$ & 0.025 & $\mathbf{1\times10^{-5}}$ & $\mathbf{0.01}$ & Large gain \\
case30\_ppcY\_A   & $3\times10^{-4}$ & 0.20  & $\mathbf{7\times10^{-5}}$ & $\mathbf{0.025}$ & Large gain \\
case33bw\_ppcY\_A & $1.5\times10^{-3}$ & 0.10 & $5\times10^{-4}$ & 0.10 & Modest gain \\
case39\_ppcY\_A   & $5\times10^{-3}$ & 0.10  & $\mathbf{5\times10^{-5}}$ & $\mathbf{0.07}$ & Large gain \\
case57\_ppcY\_A   & $3\times10^{-3}$ & 0.20  & $3\times10^{-3}$ & 0.30 & No gain \\
case89pegase      & $3\times10^{-3}$ & 0.30  & $2.5\times10^{-3}$ & 0.27 & Marginal \\
case118\_ppcY\_A  & $4\times10^{-4}$ & 0.09  & $\mathbf{6\times10^{-5}}$ & $\mathbf{0.08}$ & Large gain \\
case145\_ppcY\_A  & $0.010$ & 6.0   & $3\times10^{-3}$ & 6.0  & Modest gain \\
LVN Heo1         & --- & $\sim91$ (diverged) & \multicolumn{2}{c}{all steps rejected$^\dagger$} & No learning \\
\bottomrule
\multicolumn{6}{l}{${}^\dagger$ Armijo rejected every candidate step at every batch; model received no gradient update.}\\
\end{tabular}
\end{table}

For most cases, Armijo produces a significant accuracy improvement by
preventing divergent updates.  The exceptions are case57 and case89pegase,
which appear to be architecturally undersized at ($K=40$, $d_{\mathrm{hi}}=16$):
both cases plateau regardless of the step-size policy.  case89pegase is
notably dense ($N_{\mathrm{branch}}=210$ for 89 buses), which may
require a larger model.

\subsection{Representative Learning Curves}

Figures~\ref{fig:armijo_24}--\ref{fig:armijo_145} show the validation RMSE
curves ($|V|$ and $\theta$) for four representative cases, contrasting
Armijo=False and Armijo=True.

\begin{figure}[h]
  \centering
  \begin{subfigure}[b]{0.48\textwidth}
    \includegraphics[width=\textwidth]{case24_ieee_rts_armijoFalse_rmse}
    \caption{case24 -- Armijo=False}
  \end{subfigure}
  \hfill
  \begin{subfigure}[b]{0.48\textwidth}
    \includegraphics[width=\textwidth]{case24_ieee_rts_armijoTrue_rmse}
    \caption{case24 -- Armijo=True}
  \end{subfigure}
  \caption{case24\_ieee\_rts: Armijo reduces $|V|$ RMSE by $\sim60\times$
    and stabilises training.}
  \label{fig:armijo_24}
\end{figure}

\begin{figure}[h]
  \centering
  \begin{subfigure}[b]{0.48\textwidth}
    \includegraphics[width=\textwidth]{case39_ppcY_A_armijoFalse_rmse}
    \caption{case39 -- Armijo=False (diverges)}
  \end{subfigure}
  \hfill
  \begin{subfigure}[b]{0.48\textwidth}
    \includegraphics[width=\textwidth]{case39_ppcY_A_armijoTrue_rmse}
    \caption{case39 -- Armijo=True (converges)}
  \end{subfigure}
  \caption{case39\_ppcY\_A: without Armijo the updates overshoot and
    diverge; Armijo restores stable descent.}
  \label{fig:armijo_39}
\end{figure}

\begin{figure}[h]
  \centering
  \begin{subfigure}[b]{0.48\textwidth}
    \includegraphics[width=\textwidth]{case118_ppcY_A_armijoFalse_rmse}
    \caption{case118 -- Armijo=False}
  \end{subfigure}
  \hfill
  \begin{subfigure}[b]{0.48\textwidth}
    \includegraphics[width=\textwidth]{case118_ppcY_A_armijoTrue_rmse}
    \caption{case118 -- Armijo=True}
  \end{subfigure}
  \caption{case118\_ppcY\_A: Armijo reduces $|V|$ RMSE by $\sim7\times$
    and produces a sawtooth descent pattern consistent with cosine-restart
    scheduling.}
  \label{fig:armijo_118}
\end{figure}

\begin{figure}[h]
  \centering
  \begin{subfigure}[b]{0.48\textwidth}
    \includegraphics[width=\textwidth]{case145_ppcY_A_armijoFalse_rmse}
    \caption{case145 -- Armijo=False (oscillates)}
  \end{subfigure}
  \hfill
  \begin{subfigure}[b]{0.48\textwidth}
    \includegraphics[width=\textwidth]{case145_ppcY_A_armijoTrue_rmse}
    \caption{case145 -- Armijo=True (smoother)}
  \end{subfigure}
  \caption{case145\_ppcY\_A (small config, $d_{\mathrm{hi}}=16$, \texttt{numattn}$=1$):
    Armijo reduces $|V|$ RMSE by $\sim3\times$ but $\theta$ RMSE stays
    near $6^\circ$, indicating the model is undersized for this case.}
  \label{fig:armijo_145}
\end{figure}

%==============================================================================
\section{case145 Architecture Sweep}
%==============================================================================

\subsection{What is the Pareto Front?}

In a multi-objective setting, a run is \textbf{Pareto-optimal} (or
Pareto-dominant) if no other run is \emph{strictly better on every
objective simultaneously}.  Here the two objectives are:
\begin{enumerate}[noitemsep]
  \item Minimise $|V|$ RMSE (voltage magnitude accuracy).
  \item Minimise $\theta$ RMSE (angle accuracy).
\end{enumerate}
A run is on the Pareto front if you cannot improve one metric without
making the other worse.  Pareto-optimal runs represent the
\emph{accuracy--accuracy} trade-off frontier and are the natural
candidates for deployment depending on which metric the operator
prioritises.

\subsection{Part~A: Launcher Log Sweep}

All runs: $K=40$, 100\,epochs, Armijo$=\mathrm{False}$.
Sorted by $\theta$ RMSE (best first).
Runs with $d_{\mathrm{hi}}=64$ diverged catastrophically (physics loss
$\sim10^5$) and are listed separately.  The OOM run is excluded.

\begin{table}[h]
\centering
\scriptsize
\caption{case145 architecture sweep (launcher log sweep).
  $\checkmark$ = Pareto-optimal. $H$=heads, $L$=attn layers.}
\label{tab:case145_sweep}
\setlength{\tabcolsep}{3pt}
\begin{tabular}{clrrrrrrrrl}
\toprule
Rank & Label & $H$ & $d_{\mathrm{hi}}$ & $L$ & Params
  & $\mathcal{L}_{\mathrm{phys}}$ & $|V|$ RMSE & $\theta$ (${}^\circ$)
  & $\Delta P^\infty$ & Pareto \\
\midrule
1  & 03\_nhead8\_dhi24\_attn8   & 8 & 24 & 8  & 65,896  & 0.426 & $1.767\times10^{-3}$ & \textbf{1.770} & 2.750 & $\checkmark$ \\
2  & 09\_nhead8\_dhi32\_attn10  & 8 & 32 & 10 & 137,288 & 0.498 & $2.758\times10^{-3}$ & 2.440 & \textbf{0.320} & \\
3  & 04\_nhead8\_dhi24\_attn10  & 8 & 24 & 10 & 80,816  & 0.461 & $4.140\times10^{-3}$ & 2.776 & 0.457 & \\
4  & 08\_nhead8\_dhi32\_attn8   & 8 & 32 & 8  & 111,472 & 0.458 & $6.323\times10^{-3}$ & 2.824 & 0.990 & \\
5  & 16\_nhead12\_dhi24\_attn4  & 12& 24 & 4  & 36,360  & 0.503 & $\mathbf{1.213\times10^{-3}}$ & 3.973 & 3.374 & $\checkmark$ \\
6  & 17\_nhead12\_dhi24\_attn6  & 12& 24 & 6  & 51,432  & 0.461 & $5.496\times10^{-3}$ & 4.012 & 5.520 & \\
7  & 05\_nhead8\_dhi24\_attn12  & 8 & 24 & 12 & 95,736  & 0.431 & $3.472\times10^{-2}$ & 4.075 & 1.778 & \\
8  & 18\_nhead12\_dhi24\_attn8  & 12& 24 & 8  & 66,504  & 0.836 & $2.324\times10^{-2}$ & 4.794 & 1.114 & \\
9  & 01\_nhead8\_dhi24\_attn4   & 8 & 24 & 4  & 36,056  & 0.473 & $1.671\times10^{-3}$ & 4.945 & 1.148 & $\checkmark$ \\
10 & 06\_nhead8\_dhi32\_attn4   & 8 & 32 & 4  & 59,840  & 0.493 & $2.622\times10^{-2}$ & 5.079 & 1.350 & \\
11 & 02\_nhead8\_dhi24\_attn6   & 8 & 24 & 6  & 50,976  & 0.446 & $4.395\times10^{-3}$ & 5.411 & 4.612 & \\
12 & 07\_nhead8\_dhi32\_attn6   & 8 & 32 & 6  & 85,656  & 0.513 & $3.259\times10^{-3}$ & 6.354 & 9.195 & \\
\midrule
OOM  & 10\_nhead8\_dhi32\_attn12 & 8 & 32 & 12 & 163,104 & \multicolumn{5}{c}{out of memory} \\
FAIL & 11\_nhead8\_dhi64\_attn4  & 8 & 64 & 4 & 216,416 & $1.51\times10^5$ & 0.246 & $39.6^\circ$ & 490 & diverged \\
FAIL & 12\_nhead8\_dhi64\_attn6  & 8 & 64 & 6 & 316,536 & $1.66\times10^5$ & 0.241 & $88.0^\circ$ & 729 & diverged \\
\midrule
--- & runs 19--45 & \multicolumn{9}{c}{nhead=12 ($d_{\mathrm{hi}}\ge48$), nhead=16 -- incomplete} \\
\bottomrule
\end{tabular}
\end{table}

The three Pareto-optimal runs are:
\begin{itemize}[noitemsep]
  \item \textbf{Rank 1} ($H=8$, $d_{\mathrm{hi}}=24$, $L=8$, 65k params):
    best $\theta$ at $1.77^\circ$.
  \item \textbf{Rank 5} ($H=12$, $d_{\mathrm{hi}}=24$, $L=4$, 36k params):
    best $|V|$ RMSE at $1.21\times10^{-3}$, smallest Pareto model.
  \item \textbf{Rank 9} ($H=8$, $d_{\mathrm{hi}}=24$, $L=4$, 36k params):
    good $|V|$ RMSE ($1.67\times10^{-3}$) at minimal cost; also 36k params
    but worse than Rank 5 on $|V|$.
\end{itemize}

\begin{figure}[h]
  \centering
  \includegraphics[width=0.85\textwidth]{case145_sweep_metrics}
  \caption{case145 launcher sweep: test metrics by run
    ($K=40$, Armijo=False).  Pareto-optimal runs are the best in the
    lower-left of the $|V|$--$\theta$ plane.}
  \label{fig:sweep}
\end{figure}

\subsection{Part~B: Earlier Plot-Based Runs (small configs, numattn $= 1$--$2$)}

These runs predate the full sweep and used smaller $d_{\mathrm{hi}}$ and
fewer attention layers.

\begin{table}[h]
\centering
\caption{case145 earlier runs (from validation learning curves).}
\label{tab:case145_early}
\begin{tabular}{lrrrrrrl}
\toprule
Config & $K$ & $d$ & $d_{\mathrm{hi}}$ & \texttt{numattn} & Armijo
  & $|V|$ RMSE & $\theta$ (${}^\circ$) \\
\midrule
K40\_d4\_dhi16\_n1 & 40 & 4 & 16 & 1 & False & $\sim0.010$      & $\sim6$ \\
K40\_d4\_dhi16\_n1 & 40 & 4 & 16 & 1 & \textbf{True}  & $\sim3\times10^{-3}$ & $\sim6$ \\
K40\_d4\_dhi16\_n2 & 40 & 4 & 16 & 2 & False & $\sim2\times10^{-3}$ & $\sim5$ \\
K40\_d4\_dhi24\_n1 & 40 & 4 & 24 & 1 & False & $\sim0.010$      & $\sim5$ \\
K40\_d4\_dhi24\_n2 & 40 & 4 & 24 & 2 & False & $\sim8\times10^{-3}$ & $\sim5$ \\
K40\_d6\_dhi16\_n1 & 40 & 6 & 16 & 1 & False & $\sim0.030$      & $\sim6$ \\
K40\_d6\_dhi24\_n1 & 40 & 6 & 24 & 1 & False & $\sim0.015$      & $\sim5$ \\
K48\_d4\_dhi16\_n1 & 48 & 4 & 16 & 1 & False & $\sim3\times10^{-3}$ & $\sim5$ \\
K48\_d4\_dhi24\_n2 & 48 & 4 & 24 & 2 & False & $\sim3\times10^{-3}$ & $\sim5$ \\
K48\_d6\_dhi24\_n2 & 48 & 6 & 24 & 2 & False & $\sim0.010$      & $\sim10$ \\
\bottomrule
\end{tabular}
\end{table}

\begin{figure}[h]
  \centering
  \includegraphics[width=0.75\textwidth]{case145_incremental}
  \caption{case145 incremental run summary (earlier small-config sweep).
    All configs plateau at $\theta \approx 5$--$6^\circ$, well above the
    best from the full capacity sweep ($1.77^\circ$).}
  \label{fig:incremental}
\end{figure}

\subsection{Key Findings}

\begin{enumerate}[noitemsep]
  \item \textbf{Best overall:} $H{=}8$, $d_{\mathrm{hi}}{=}24$, $L{=}8$
    (65k params) achieves $\theta = 1.77^\circ$, $|V| = 1.77\times10^{-3}$.
  \item \textbf{Best $|V|$ RMSE:} $H{=}12$, $d_{\mathrm{hi}}{=}24$, $L{=}4$
    (36k params) achieves $|V| = 1.21\times10^{-3}$, but $\theta = 3.97^\circ$.
  \item \textbf{Deeper is not always better:} for $d_{\mathrm{hi}}=24$,
    $L{=}8$ outperforms both $L{=}10$ and $L{=}12$ on angle accuracy.
    $L{=}12$ has 3$\times$ the angle RMSE of $L{=}8$.
  \item \textbf{$d_{\mathrm{hi}}=64$ diverged completely:} physics loss jumped
    to $\sim10^5$, $\theta$ errors of $40$--$88^\circ$.  The likely cause
    is that the large hidden dimension magnifies gradient variance at
    initialisation (all heads initialised to zero), making the early
    training unstable under the physics loss only.
  \item \textbf{Small configs plateau at $\theta \approx 5$--$6^\circ$:}
    the earlier runs with $d_{\mathrm{hi}} \le 16$ and \texttt{numattn}$\le2$
    all converge to this plateau, confirming that architectural capacity
    (not just steps $K$) is the bottleneck for case145.
  \item \textbf{Many sweep runs still incomplete:} nhead$=12$ with
    $d_{\mathrm{hi}}\ge48$ and all nhead$=16$ runs are pending; better
    configurations may yet be found.
\end{enumerate}

%==============================================================================
\section{case300 Results (PEGASE 300-bus, A100)}
%==============================================================================

\subsection{Setup and Baseline}

The DC-initialised baseline for case300 (computed from 36,000 scenarios,
300 buses each) is:
\[
  |V|_{\mathrm{RMSE}}^{\mathrm{base}} = 4.82\times10^{-2}\;\text{p.u.},
  \qquad
  \theta_{\mathrm{RMSE}}^{\mathrm{base}} = 19.58^\circ
\]
This is comparable to case145 ($19.58^\circ$ vs $18.64^\circ$) because
PEGASE networks of similar voltage-level complexity share similar angle
spreads at operating conditions.

case300 (\texttt{case300\_ppcY\_A}) is the largest case tested so far:
300 buses and 411 branches.  All six runs were submitted to the \texttt{alex}
cluster on a 40\,GB A100 GPU under a 24\,h walltime limit
(sweep \texttt{case300\_top3\_armijo\_compare\_20260324\_210253}).
The configurations are the three best-performing architectures from the
case145 sweep, each tested with and without Armijo line search.

\subsection{Run-by-Run Results}

\begin{table}[h]
\centering
\small
\caption{case300 results. $K=40$, $d=4$, $H=8$.
  ``Last valid'' means the last epoch logged before timeout.
  The only fully completed run is job \texttt{3492156}.}
\label{tab:case300}
\setlength{\tabcolsep}{4pt}
\begin{tabular}{llrrrlll}
\toprule
Job & Config & $d_{\mathrm{hi}}$ & $L$ & Epoch & $|V|$ RMSE (p.u.) & $\theta$ (${}^\circ$) & Status \\
\midrule
3492155 & noarmijo & 24 & 8  & 94  & $2.171\times10^{-2}$ & 1.806 & TIMEOUT \\
\textbf{3492156} & \textbf{armijo}   & \textbf{24} & \textbf{8}  & \textbf{100} & $\mathbf{9.800\times10^{-3}}$ & \textbf{1.774} & \textbf{COMPLETE}$^\dagger$ \\
3492157 & noarmijo & 32 & 10 & 80  & $1.993\times10^{-2}$ & 3.211 & TIMEOUT \\
3492158 & armijo   & 32 & 10 & 91  & $1.196\times10^{-2}$ & 4.101 & TIMEOUT \\
3492159 & noarmijo & 24 & 10 & 74  & $2.067\times10^{-2}$ & 3.423 & TIMEOUT \\
3492160 & armijo   & 24 & 10 & 94  & $1.468\times10^{-2}$ & 1.825 & TIMEOUT \\
\bottomrule
\multicolumn{7}{l}{${}^\dagger$ Slurm reports \texttt{FAILED 120} (walltime signal), but the run finished all 100 epochs}\\
\multicolumn{7}{l}{\phantom{${}^\dagger$} and printed final test metrics before exit; the result is usable.}\\
\end{tabular}
\end{table}

\subsection{Armijo Effect on case300}

On case300 Armijo consistently helps, which is even clearer than on case145:

\begin{center}
\begin{tabular}{lrrrrr}
\toprule
Config & Armijo & Epoch & $|V|$ RMSE & $\theta$ (${}^\circ$) & $\Delta\theta$ \\
\midrule
$d_{\mathrm{hi}}=24$, $L=8$  & No  & 94  & $2.171\times10^{-2}$ & 1.806 & \multirow{2}{*}{$-0.032$} \\
$d_{\mathrm{hi}}=24$, $L=8$  & Yes & 100 & $9.800\times10^{-3}$ & \textbf{1.774} & \\
\midrule
$d_{\mathrm{hi}}=24$, $L=10$ & No  & 74  & $2.067\times10^{-2}$ & 3.423 & \multirow{2}{*}{$-1.598$} \\
$d_{\mathrm{hi}}=24$, $L=10$ & Yes & 94  & $1.468\times10^{-2}$ & \textbf{1.825} & \\
\midrule
$d_{\mathrm{hi}}=32$, $L=10$ & No  & 80  & $1.993\times10^{-2}$ & 3.211 & \multirow{2}{*}{$-0.890$} \\
$d_{\mathrm{hi}}=32$, $L=10$ & Yes & 91  & $1.196\times10^{-2}$ & \textbf{4.101}$^\ddagger$ & \\
\bottomrule
\multicolumn{6}{l}{${}^\ddagger$ Job \texttt{3492158} showed better $\theta$ around epoch 79; later epochs degraded.}
\end{tabular}
\end{center}

The Armijo run with $d_{\mathrm{hi}}=24$, $L=8$ reduces $|V|$ RMSE by
$\sim2.2\times$ compared to its no-Armijo counterpart (both at nearly the
same epoch count), and achieves the best $\theta$ of any case300 run.

\subsection{Throughput Bottleneck Analysis}

Five of six jobs timed out under the 24\,h walltime.  Profiling of job
\texttt{3492160} reveals that the limiting factor is \emph{not} GPU memory or
compute capacity:

\begin{itemize}[noitemsep]
  \item \textbf{GPU memory:} $\sim7$\,GB used out of 40\,GB available on
    the A100.  The GPU was $\sim83\%$ idle on memory.
  \item \textbf{GPU utilisation:} $35$--$50\%$ compute utilisation with
    frequent dips.  The GPU was not saturated.
  \item \textbf{Host RAM:} \texttt{sacct} reports $\sim86$--$88$\,GB RSS
    on the host CPU.  Heavy CPU-side overhead from data loading,
    graph construction, or Python dispatch is the likely bottleneck.
\end{itemize}

This is a \emph{throughput bottleneck}, not a capacity bottleneck.
The GPU waits for the CPU.  Practical remedies:
\begin{enumerate}[noitemsep]
  \item Increase batch size to better saturate the A100.
  \item Profile and reduce CPU-side graph-construction or collation cost.
  \item Request longer walltime ($\ge 48$\,h) to allow full training on the
    current configuration.
\end{enumerate}

\subsection{Key Findings}

\begin{enumerate}[noitemsep]
  \item \textbf{Best result:} $d_{\mathrm{hi}}=24$, $L=8$, Armijo=True ---
    $|V|=9.80\times10^{-3}$\,p.u., $\theta=1.774^\circ$ at epoch 100
    (the only fully completed run).
  \item \textbf{Second best:} $d_{\mathrm{hi}}=24$, $L=10$, Armijo=True ---
    $\theta=1.825^\circ$ at epoch 94 (timed out; likely would improve
    further with full 100 epochs).
  \item \textbf{Armijo is essential at this scale:} across all three
    architecture pairs, Armijo reduces $\theta$ RMSE by $0.03$--$1.6^\circ$.
  \item \textbf{$d_{\mathrm{hi}}=24$, $L=8$ remains the best architecture}
    on case300, consistent with the finding from case145.
  \item \textbf{Job \texttt{3492158} ($d_{\mathrm{hi}}=32$, Armijo) degraded
    after epoch 79:} the best checkpoint is likely better than the final
    logged value; saving and evaluating mid-training checkpoints is
    recommended for long jobs.
  \item \textbf{Runtime, not memory, is the constraint on alex:}
    increasing walltime or improving CPU throughput is the immediate next
    step before scaling further.
\end{enumerate}

%==============================================================================
\section{LVN Experiments: An Open Problem}
\label{sec:lvn}
%==============================================================================

\subsection{Overview}

\textbf{LVN} stands for \textit{LEW Verteilnetz GmbH}, the electricity
distribution subsidiary of Lechwerke AG (LEW), a major regional energy company
in Bavaria, Germany.  The LVN Heo1 case is a real-world distribution network
operated by LEW Verteilnetz GmbH, and presents a fundamentally harder
power-flow problem than the HV/MV PEGASE benchmark cases above.  Three
experimental series were conducted, progressively diagnosing the failure modes.
The conclusion is that \textbf{none of the LVN experiments produced a working
physics-only solver} under the current architecture.

\subsection{Series 1 \& 2: DC-Informed Initialization
  (\texttt{lvn\_case300\_informed})}

The first two folders (\texttt{20260422\_080112} and
\texttt{20260422\_083606}) used the DC power-flow solution as $V_{\rm
start}$ (``informed'' initialization).  These runs exposed two distinct
failure modes.

\paragraph{Armijo runs: physics loss always detached.}
All Armijo configurations produced an unbroken stream of
\texttt{[warn] physics loss detached} messages from the very first batch.
This means the Armijo line search rejected \emph{every} candidate step at
every batch: the $\ell^\infty$-norm of the power mismatch did not decrease
for any step size in the candidate set, so the model never received a
gradient.  The network made zero progress.

\paragraph{No-Armijo runs: immediate divergence to large angle errors.}
Without the line search, updates were accepted unconditionally.  However,
angle RMSE sat at $34$--$52^\circ$ throughout training, and the physics
loss was enormous ($\sim10^{15}$--$10^{16}$).  The model was updating but
not descending the residual.

\paragraph{The DC-init artifact.}
The second folder revealed a subtler issue: with DC-initialized
$V_{\rm start}$, the Armijo runs logged epoch~0 validation metrics of
$|V|_{\rm RMSE} = 4.48\times10^{-2}$\,p.u.\ and $\theta = 0.24^\circ$,
with $\Delta P^\infty = 6\times10^3$\,p.u.  This looks deceptively good,
but the $0.24^\circ$ angle error was already present in $V_{\rm start}$
before the model did anything.  The DC solution was already very close to
the AC-NR solution for this particular dataset, so there was almost nothing
to learn---and the residuals were still $\sim6000$\,p.u., indicating the
network had not reduced the mismatch at all.

\subsection{Series 3: Manual Flat Initialization with K-Hop
  (\texttt{lvn\_manual\_flat\_khop\_ablation})}

To remove the DC-init artifact and test a genuinely hard starting point,
the third series used manual flat initialization ($|V|=1.0$\,p.u.,
$\theta=0^\circ$) and added K-hop Gaussian diffusion augmentation (see
Section~\ref{} in the khop document).  Twelve configurations were tested,
varying $K_{\rm hop} \in \{2, 3, 4\}$, $\sigma \in \{1.0, 1.5\}$, loss
form (\texttt{mse} / \texttt{huber}+setpoint normalization), and final
residual penalty.

All runs timed out between epoch 8 and epoch 14.
Table~\ref{tab:lvn_khop} shows the key metrics.

\begin{table}[h]
\centering
\scriptsize
\caption{LVN K-hop ablation (manual flat init, partial runs).
  All $\theta$ values are in degrees.
  Residuals in p.u.\ --- note the $\sim10^6$ scale.}
\label{tab:lvn_khop}
\setlength{\tabcolsep}{3pt}
\begin{tabular}{lrrllrrrr}
\toprule
Config & $K_{\rm hop}$ & $\sigma$ & Loss & Norm
  & Ep. & Best $\theta$ ($^\circ$) & Mean $|\Delta P|$ (pu) & Mean $|\Delta Q|$ (pu) \\
\midrule
khop2 sig1.0 mse        & 2 & 1.0 & mse   & none     & 15 & 91.2  & $3.4\times10^6$ & $6.5\times10^6$ \\
khop3 sig1.0 mse        & 3 & 1.0 & mse   & none     & 12 & 92.3  & $3.6\times10^6$ & $5.2\times10^6$ \\
khop4 sig1.0 mse        & 4 & 1.0 & mse   & none     & 10 & 95.1  & $3.1\times10^6$ & $6.0\times10^6$ \\
khop2 sig1.0 huber+sp   & 2 & 1.0 & huber & setpoint & 15 & 98.6  & $4.1\times10^6$ & $5.8\times10^6$ \\
khop2 sig1.5 mse        & 2 & 1.5 & mse   & none     & 14 & 99.1  & $4.0\times10^6$ & $6.6\times10^6$ \\
khop4 sig1.0 huber+sp   & 4 & 1.0 & huber & setpoint & 9  & 99.8  & $4.0\times10^6$ & $6.0\times10^6$ \\
khop3 sig1.0 huber+sp   & 3 & 1.0 & huber & setpoint & 12 & 100.6 & $3.6\times10^6$ & $6.3\times10^6$ \\
khop3 sig1.5 mse        & 3 & 1.5 & mse   & none     & 11 & 102.0 & $3.7\times10^6$ & $5.9\times10^6$ \\
khop4 sig1.5 mse        & 4 & 1.5 & mse   & none     & 9  & 102.7 & $3.7\times10^6$ & $5.7\times10^6$ \\
khop3 sig1.5 huber+sp   & 3 & 1.5 & huber & setpoint & 11 & 104.0 & $3.3\times10^6$ & $6.1\times10^6$ \\
khop2 sig1.5 huber+sp   & 2 & 1.5 & huber & setpoint & 14 & 105.2 & $3.4\times10^6$ & $6.0\times10^6$ \\
khop4 sig1.5 huber+sp   & 4 & 1.5 & huber & setpoint & 9  & 105.2 & $3.2\times10^6$ & $6.4\times10^6$ \\
\bottomrule
\end{tabular}
\end{table}

The best configuration (\texttt{khop2}, $\sigma{=}1.0$, raw MSE) achieved
a best validation $\theta = 91.2^\circ$---essentially no improvement over
the flat-start baseline of $\sim90^\circ$ (LVN Heo1 uses a genuine flat start, not DC init).  Mean $|\Delta P|$ across all
runs was $3$--$4\times10^6$\,p.u., and the fraction of buses below
convergence tolerance was $0\%$.  K-hop augmentation did not unlock any
learning.

\begin{figure}[h]
  \centering
  \begin{subfigure}[b]{0.48\textwidth}
    \includegraphics[width=\textwidth]{lvn_khop_rmse_curves}
    \caption{Validation RMSE curves (all K-hop configs)}
  \end{subfigure}
  \hfill
  \begin{subfigure}[b]{0.48\textwidth}
    \includegraphics[width=\textwidth]{lvn_khop_residual_curves}
    \caption{Mean $|\Delta P|$, $|\Delta Q|$ residual curves}
  \end{subfigure}
  \caption{LVN K-hop ablation: validation RMSE (left) and mean power
    residuals (right) across all 12 configurations.  No run shows
    meaningful descent.  The $\sim10^6$\,p.u.\ residuals indicate the
    model is not approaching AC-PF feasibility.}
  \label{fig:lvn_khop}
\end{figure}

\begin{figure}[h]
  \centering
  \includegraphics[width=0.6\textwidth]{lvn_khop_ranking}
  \caption{Best-angle ranking across K-hop configs.  All best angles
    cluster near $91$--$107^\circ$, confirming none approach a useful
    solution.}
  \label{fig:lvn_ranking}
\end{figure}

\subsection{Lessons Learned}

\begin{enumerate}[noitemsep]
  \item \textbf{Ybus bug was real but not the whole problem.}
    LVN uses asymmetric branch parameters.  The old reconstructed Ybus was
    wrong.  After fixing directed branch admittance, the power-flow
    residuals became consistent with $V_{\rm newton}$.  Training still did
    not learn.

  \item \textbf{DC init was misleading.}
    The $0.24^\circ$ angle RMSE in the informed runs was already in
    $V_{\rm start}$, not earned by the model.  The flat-start series gives
    the honest benchmark.

  \item \textbf{Armijo failed for LVN.}
    The line search rejected all steps in every batch for every Armijo
    run.  The proposed update always increased the mismatch $\ell^\infty$-norm.
    This indicates the GNN is not producing reliable descent directions,
    not merely choosing wrong step sizes.

  \item \textbf{Capacity did not transfer.}
    The case145-best settings ($d_{\rm hi}=24$, $H=8$, $L=8$, $K=40$) do
    not carry over to LVN.  The problem is not ``use a bigger model.''

  \item \textbf{Loss modifications did not fix it.}
    Raw MSE, Huber, setpoint normalization, graph normalization, and final
    residual penalty all failed to produce convergence.  The loss landscape
    is dominated by extreme residuals and the gradient signal is
    uninformative.

  \item \textbf{K-hop did not help.}
    Multi-hop electrical context ($K_{\rm hop}=2$--$4$) added no useful
    information beyond 1-hop attention.  The bottleneck is not receptive
    field.
\end{enumerate}

\subsection{Interpretation and Next Step}

The evidence points to a learning-dynamics failure, not a capacity or
receptive-field failure.  The model is not producing stable correction
directions for LVN under the current physics-only objective.  The unrolled
updates appear to be exploring the loss landscape rather than descending it.

The recommended next direction is to shift the update mechanism from
free-form voltage corrections to a \textbf{physically preconditioned
update}---one that normalizes or constrains the per-bus update by bus type,
apparent power, or a residual-conditioned damping.  This would give the
gradient a more informative and consistent structure than the raw
AC-residual MSE currently provides.  Increasing model size should come
after this structural change, not before.

%==============================================================================
\section{LVN Heo1 Grid: First Successful Run}
%==============================================================================

\subsection{Dataset and Setup}

The Heo1 grid is a real-world LVN case with multi-voltage topology spanning
$3$--$380\,\mathrm{kV}$.  Training used the flat-start parquet:

\begin{center}
\small\ttfamily\raggedright
LVN\_snapshot\_envelope\_manual\_flat\_coupled\_ppcY\_manual\_flat\_%
siNR\_36000\_NR\_branchrows\_directSI.parquet
\end{center}

36{,}000 scenarios, train/val/test split $1/3$ each.

\subsection{Bugs Found and Fixed Before This Run}

Three correctness issues were identified and fixed before the Heo1
experiment, which explain why all prior LVN runs failed to learn.

\paragraph{Issue 1: Bus-type mapping was inverted.}
The bus-type integers for Slack, PQ, and PV buses were mapped in the wrong
order in the LVN dataset pipeline.  This meant the model zeroed $\Delta P$
and $\Delta Q$ at the wrong buses at every step---the physics-informed
features were corrupted from the start.

\paragraph{Issue 2: LVN multi-voltage outlier filter.}
LVN spans $3$--$380\,\mathrm{kV}$.  The global $|V|$ outlier filter
applied during preprocessing used thresholds calibrated for single-voltage
HV cases.  In the multi-voltage Heo1 grid, low-voltage buses ($3\,\mathrm{kV}$)
had per-unit magnitudes far outside the HV range, causing many valid
scenarios to be incorrectly filtered or their voltage features to be
clipped.

\paragraph{Issue 3: Pure PINN diverged on LVN.}
Physics-only training diverged immediately: validation RMSE jumped from
$0.073$ at epoch~0 to $0.142$ at epoch~1 and continued growing.  LVN
Jacobians are stiffer than HV/MV cases because of the wide voltage range
and mixed transformer ratios.  The raw AC-residual loss landscape is too
sharp for the GNN to descend without a supervised anchor.

\subsection{Training Strategy: Combined Loss}

To stabilize training after the pure-PINN failure, a combined loss was
used:

\begin{equation}
  \mathcal{L}
  =
  \mathcal{L}_{\rm phys}
  +
  w \cdot \mathcal{L}_{\rm MSE}
  \label{eq:combined}
\end{equation}

with $w = 10$.  The supervised MSE term ($\mathcal{L}_{\rm MSE}$, comparing
predicted voltages to the NR reference $V_{\rm newton}$) acts as a
stabilizing anchor that keeps the network near the correct solution while
the physics loss provides residual-based gradient information.

\subsection{Results}

\subsubsection{Test Set Performance (best.ckpt at Epoch 24)}

\begin{table}[h]
\centering
\caption{LVN Heo1 test results from best checkpoint (epoch~24).}
\label{tab:heo1_test}
\begin{tabular}{lc}
\toprule
Metric & Value \\
\midrule
Test RMSE (total)    & $1.98\times10^{-2}$ \\
Test $|V|$ RMSE      & $2.43\times10^{-2}$\,p.u. \\
Test $\theta$ RMSE   & $0.794^\circ$ \\
Test physics loss    & $3.671\times10^{5}$ \\
\bottomrule
\end{tabular}
\end{table}

This is the first LVN run to achieve sub-degree angle RMSE.

\subsubsection{Pure MSE vs.\ Combined Loss}

\begin{table}[h]
\centering
\caption{Comparison of training strategies on LVN Heo1.}
\label{tab:heo1_compare}
\begin{tabular}{lccc}
\toprule
Strategy & RMSE & $\theta$ RMSE & Physics loss \\
\midrule
Pure MSE ($w=0$)           & $0.020$ & $0.79^\circ$ & huge (not minimized) \\
Combined phys $+$ $10\cdot$MSE & ${\sim}0.040$ & ${\sim}2.1^\circ$ & ${\sim}0.9$ \\
\bottomrule
\end{tabular}
\end{table}

The comparison reveals a fundamental tension:

\begin{itemize}[noitemsep]
  \item \textbf{Pure MSE} achieves the best supervised accuracy ($\theta =
    0.79^\circ$) because the entire gradient budget is spent imitating
    $V_{\rm newton}$.  However, the physics loss remains enormous---the
    model has learned to copy NR output without learning to satisfy AC-PF
    equations independently.

  \item \textbf{Combined loss} achieves a lower physics loss (${\sim}0.9$
    vs.\ huge) at the cost of higher supervised error ($\theta \approx
    2.1^\circ$).  The model makes partial progress toward AC-PF
    feasibility, but the MSE anchor prevents it from fully committing to
    residual descent.
\end{itemize}

\subsection{Interpretation}

The physics loss of $3.671\times10^5$ in the best-checkpoint test result
confirms the model has not yet solved LVN from a physics-only standpoint,
but the $0.794^\circ$ angle RMSE demonstrates that the combined-loss
approach stabilizes training where pure PINN failed.  The three bugs
(bus-type inversion, outlier filter, pure-PINN divergence) were the primary
blockers for all prior LVN experiments---not model capacity or receptive
field, as previously suspected.

\paragraph{Next steps.}
With the bugs fixed and a working training loop, the productive direction
is to reduce $w$ gradually (annealing the MSE weight toward zero) so the
model transitions from supervised imitation to physics-driven correction.
Alternatively, the combined loss can be kept but the physics weight
increased to push the physics loss below an operational threshold.

\end{document}
